% document formatting
\documentclass[10pt]{article}
\usepackage[utf8]{inputenc}
\usepackage[left=1in,right=1in,top=1in,bottom=1in]{geometry}
\usepackage[T1]{fontenc}
\usepackage{xcolor}

% math symbols, etc.
\usepackage{amsmath, amsfonts, amssymb, amsthm}

% lists
\usepackage{enumerate}

% images
\usepackage{graphicx} % for images
\usepackage{tikz} 

% code blocks
% \usepackage{minted, listings} 
\usepackage[colorlinks=true, urlcolor=blue]{hyperref}


% verbatim greek
\usepackage{alphabeta}

\graphicspath{{./assets/images/Week 8}}

\title{02-719 Week 8 \\ \large{Genomics and Epigenetics of the Brain}}
  
\author{Aidan Jan}

\date{\today}

\begin{document}
\maketitle
\subsection*{Recap: Convergent Evolution}
\begin{itemize}
    \item The same phenotype evolves \textbf{independently} in separate lineages
    \item Each origin is a natural replicate of the same experiment
    \begin{itemize}
        \item If a gene matters for the trait, it may respond the same way each time
    \end{itemize}
    \item The question: Does a gene's \textbf{rate of evolution} shift on the branches where are trait evolved?
\end{itemize}

\subsection*{Species are not independent data points}
Naive approach:
\begin{itemize}
    \item Correlate a gene's sequence with the trait across species
    \item Closely related species share both, so one origin can look like many
    \item $\rightarrow$ inflated significance
\end{itemize}
Phylogenetic approach:
\begin{itemize}
    \item Work on branches of the tree, not species
    \item Count each independent origin once
    \item $\rightarrow$ signal must repeat across lineages
\end{itemize}

\section*{RER Coverage}
RER = Relative evolutionary rates.  \textbf{RERcoverage} is an R workflow that does the following analysis:
\begin{itemize}
    \item Inputs: gene trees with a fixed topology
    \item Relative evolutionary rates (RERs)
    \item Encoding the phenotype (binary vs continuous traits)
    \item Testing for association
    \item From gene lists to biology
\end{itemize}

\subsection*{RERcoverage Workflow}
\begin{enumerate}
    \item Alignments: one per gene
    \item Gene trees: fixed topology, gene-specific branch lengths
    \item RERs: relative rate on every branch
    \item Phenotype paths: trait mapped onto branches
    \item Correlation: RER vs. trait, per gene
    \item Enrichment: Pathways / gene ontology
\end{enumerate}
Each gene gets one test.  The output is a ranked list of genes whose rates shift with the trait across the tree.

\subsection*{Step 1: Inputs, one tree per gene}
\begin{itemize}
    \item Start from a multiple alignment of orthologs for each gene
    \item Fix the \textbf{species tree topology}; estimate only the branch lengths
    \begin{itemize}
        \item Branch length = substitutions per site on that branch
        \item Species missing the gene are pruned for that gene
    \end{itemize}
    \item readTrees() combines all genes
    \begin{itemize}
        \item Builds a \textbf{master tree}: average branch lengths across genes
    \end{itemize}
\end{itemize}




\end{document}